\documentclass[12pt,a4paper]{article}

% ============================================================================
%  Complete mathematical derivation for the PDM (Patch-Diffusion) UAD pipeline
%  Self-contained; can also be \input{} into the main PDM report (delete the
%  preamble + \begin/\end{document} when including).
% ============================================================================

\usepackage[utf8]{inputenc}
\usepackage[a4paper,left=22mm,right=22mm,top=22mm,bottom=24mm]{geometry}
\usepackage{amsmath,amssymb,amsthm,mathtools}
\usepackage{bm}
\usepackage{graphicx}
\usepackage{booktabs}
\usepackage{xcolor}
\usepackage{hyperref}
\hypersetup{colorlinks=true,linkcolor=blue!50!black,citecolor=green!50!black,urlcolor=red!50!black}
\usepackage{titlesec}
\usepackage{enumitem}
\usepackage{times}

\newtheorem{proposition}{Proposition}
\newcommand{\R}{\mathbb{R}}
\newcommand{\E}{\mathbb{E}}
\newcommand{\N}{\mathcal{N}}
\newcommand{\KL}{\mathrm{D_{KL}}}
\newcommand{\x}{\bm{x}}
\newcommand{\z}{\bm{z}}
\newcommand{\eps}{\bm{\epsilon}}
\newcommand{\II}{\mathbf{I}}
\newcommand{\abar}{\bar{\alpha}}
\newcommand{\given}{\,\vert\,}
\DeclareMathOperator*{\argmin}{arg\,min}
\DeclareMathOperator{\ReLU}{ReLU}
\DeclareMathOperator{\Var}{Var}

\title{\textbf{Mathematical Foundations of a Patch-Diffusion Model for\\
Unsupervised Anomaly Detection in Pediatric Brain MRI}\\[4pt]
\large A complete derivation: from the DDPM variational bound to multi-scale
counterfactual scoring}
\author{Soham Patel \\ \small Deep Vision (Summer 2025) --- OTH Amberg-Weiden}
\date{\today}

\titleformat{\section}{\normalfont\large\bfseries}{\thesection}{1em}{}
\titleformat{\subsection}{\normalfont\normalsize\bfseries}{\thesubsection}{1em}{}

\begin{document}
\maketitle
\tableofcontents
\bigskip

\begin{abstract}
\noindent This document develops, from first principles, the complete
mathematics of the patch-diffusion model (PDM) used for unsupervised anomaly
detection (UAD) in this project. Part~\ref{part:ddpm} derives the denoising
diffusion probabilistic model (DDPM): the forward Gaussian corruption and its
closed-form marginal, the variational lower bound, the tractable forward
posterior, the reduction of the bound to a noise-prediction regression, and the
cosine noise schedule. Part~\ref{part:ddim} derives the deterministic DDIM
sampler used at inference. Part~\ref{part:pdm} specialises the theory to our
system: correlated (simplex) noise, patch-wise diffusion with Gaussian fusion,
the partial-noise reconstruction principle that makes anomalies ``heal,'' the
multi-scale modality-weighted anomaly score, the healthy-set calibration, and
the evaluation metrics. Part~\ref{part:xai} gives the mathematics of the
explainability layer (counterfactual maps, Grad-CAM for a denoiser, additive
modality/scale attribution). Part~\ref{part:numbers} explains how to read every
number the pipeline produces during training and inference.
\end{abstract}

\section*{Notation}
\addcontentsline{toc}{section}{Notation}
A clean MRI patch (or slice) is $\x_0\in\R^{C\times P\times P}$ with $C=4$ co-registered
modalities $(t1n,t1c,t2w,t2f)$, $P=96$ (patch) or $256$ (slice). The diffusion
index is $t\in\{0,\dots,T\}$, $T=1000$. $\eps$ denotes injected noise, $\eps_\theta$
the network's prediction. $\N(\bm\mu,\bm\Sigma)$ is a Gaussian; $\II$ the identity.
Per-step retention $\alpha_t=1-\beta_t$ and cumulative retention
$\abar_t=\prod_{s=1}^{t}\alpha_s$.

% ===========================================================================
\part*{Part I\quad The DDPM}
\addcontentsline{toc}{section}{Part I — The DDPM}
\label{part:ddpm}
% ===========================================================================

\section{Generative model and the forward process}
A diffusion model is a latent-variable model
$p_\theta(\x_0)=\int p_\theta(\x_{0:T})\,d\x_{1:T}$ whose latents $\x_{1:T}$ are
produced by a \emph{fixed} Markov chain that gradually destroys structure
(Sohl-Dickstein et al., 2015; Ho et al., 2020). Each forward step adds a small
amount of Gaussian noise,
\begin{equation}
q(\x_t\given \x_{t-1}) = \N\!\big(\x_t;\, \sqrt{1-\beta_t}\,\x_{t-1},\; \beta_t \II\big),
\qquad
q(\x_{1:T}\given\x_0)=\prod_{t=1}^{T} q(\x_t\given\x_{t-1}),
\label{eq:fwd-step}
\end{equation}
with a variance schedule $0<\beta_1<\dots<\beta_T<1$. The scaling $\sqrt{1-\beta_t}$
(rather than $1$) keeps the process variance-preserving, so that $\x_T$ converges
to $\N(\bm 0,\II)$ regardless of $\x_0$.

\subsection{Closed-form marginal \texorpdfstring{$q(\x_t\given\x_0)$}{q(x\_t | x\_0)}}
The single most useful property of \eqref{eq:fwd-step} is that we can sample
$\x_t$ at \emph{any} $t$ directly, without simulating the chain.

\begin{proposition}[Diffusion kernel]
With $\alpha_t=1-\beta_t$ and $\abar_t=\prod_{s=1}^t \alpha_s$,
\begin{equation}
q(\x_t\given\x_0)=\N\!\big(\x_t;\,\sqrt{\abar_t}\,\x_0,\;(1-\abar_t)\II\big)
\;\Longleftrightarrow\;
\x_t=\sqrt{\abar_t}\,\x_0+\sqrt{1-\abar_t}\,\eps,\quad \eps\sim\N(\bm0,\II).
\label{eq:kernel}
\end{equation}
\end{proposition}
\begin{proof}
Write one step as $\x_t=\sqrt{\alpha_t}\,\x_{t-1}+\sqrt{1-\alpha_t}\,\eps_{t-1}$ with
$\eps_{t-1}\sim\N(\bm0,\II)$ i.i.d. Unrolling once,
\[
\x_t=\sqrt{\alpha_t\alpha_{t-1}}\,\x_{t-2}
+\underbrace{\sqrt{\alpha_t(1-\alpha_{t-1})}\,\eps_{t-2}+\sqrt{1-\alpha_t}\,\eps_{t-1}}_{=: \,\bm u}.
\]
$\bm u$ is a sum of independent zero-mean Gaussians, hence Gaussian with variance
$\alpha_t(1-\alpha_{t-1})+(1-\alpha_t)=1-\alpha_t\alpha_{t-1}$. Thus
$\x_t=\sqrt{\alpha_t\alpha_{t-1}}\,\x_{t-2}+\sqrt{1-\alpha_t\alpha_{t-1}}\,\bar\eps$.
Iterating down to $\x_0$ replaces the product by $\abar_t$ and the variance by
$1-\abar_t$, giving \eqref{eq:kernel}.
\end{proof}

\noindent Equation~\eqref{eq:kernel} is what makes training cheap: for a random
$t$ we obtain a training pair $(\x_t,\eps)$ in one shot. The
\emph{signal-to-noise ratio} of the corrupted sample is
$\mathrm{SNR}(t)=\abar_t/(1-\abar_t)$, monotonically decreasing in $t$; it
governs which structures survive to step $t$ and is central to the UAD argument
in \S\ref{sec:partial}.

\section{The reverse process and the variational bound}
Since $q(\x_{t-1}\given\x_t)$ is intractable, we learn a Gaussian reverse chain
\begin{equation}
p_\theta(\x_{t-1}\given\x_t)=\N\!\big(\x_{t-1};\,\bm\mu_\theta(\x_t,t),\,\sigma_t^2\II\big),
\qquad p(\x_T)=\N(\bm0,\II).
\label{eq:rev}
\end{equation}
We fit $\theta$ by maximising a variational lower bound on $\log p_\theta(\x_0)$.

\begin{proposition}[Variational bound]
\begin{equation}
\E_q\!\big[-\log p_\theta(\x_0)\big]\le
\E_q\Big[\underbrace{\KL\!\big(q(\x_T\given\x_0)\,\|\,p(\x_T)\big)}_{L_T}
+\!\sum_{t=2}^{T}\underbrace{\KL\!\big(q(\x_{t-1}\given\x_t,\x_0)\,\|\,p_\theta(\x_{t-1}\given\x_t)\big)}_{L_{t-1}}
\underbrace{-\log p_\theta(\x_0\given\x_1)}_{L_0}\Big].
\label{eq:elbo}
\end{equation}
\end{proposition}
\begin{proof}[Sketch]
Start from $-\log p_\theta(\x_0)\le -\log p_\theta(\x_0)+\KL\big(q(\x_{1:T}\given\x_0)\|p_\theta(\x_{1:T}\given\x_0)\big)
=\E_q[-\log \frac{p_\theta(\x_{0:T})}{q(\x_{1:T}\given\x_0)}]$.
Expand both products over $t$, rewrite the forward chain with Bayes' rule as
$q(\x_t\given\x_{t-1})=q(\x_{t-1}\given\x_t,\x_0)\,q(\x_t\given\x_0)/q(\x_{t-1}\given\x_0)$,
and telescope the $q(\x_t\given\x_0)$ ratios. Grouping terms yields
\eqref{eq:elbo}. $L_T$ has no parameters; the work is in the $L_{t-1}$ terms.
\end{proof}

\subsection{The tractable forward posterior \texorpdfstring{$q(\x_{t-1}\given\x_t,\x_0)$}{q(x\_\{t-1\} | x\_t, x\_0)}}
Each $L_{t-1}$ compares the learned reverse step to the \emph{true} posterior
\emph{conditioned on $\x_0$}, which \emph{is} tractable.

\begin{proposition}
\begin{equation}
q(\x_{t-1}\given\x_t,\x_0)=\N\!\big(\x_{t-1};\,\tilde{\bm\mu}_t(\x_t,\x_0),\,\tilde\beta_t\II\big),
\end{equation}
\begin{equation}
\tilde\beta_t=\frac{1-\abar_{t-1}}{1-\abar_t}\,\beta_t,
\qquad
\tilde{\bm\mu}_t(\x_t,\x_0)=\frac{\sqrt{\abar_{t-1}}\,\beta_t}{1-\abar_t}\,\x_0
+\frac{\sqrt{\alpha_t}\,(1-\abar_{t-1})}{1-\abar_t}\,\x_t.
\label{eq:posterior}
\end{equation}
\end{proposition}
\begin{proof}
By Bayes,
$q(\x_{t-1}\given\x_t,\x_0)\propto q(\x_t\given\x_{t-1})\,q(\x_{t-1}\given\x_0)$,
a product of two Gaussians (in $\x_{t-1}$) from \eqref{eq:fwd-step} and
\eqref{eq:kernel}. Collecting the exponent,
\[
-\tfrac12\!\left[\frac{\|\x_t-\sqrt{\alpha_t}\x_{t-1}\|^2}{\beta_t}
+\frac{\|\x_{t-1}-\sqrt{\abar_{t-1}}\x_0\|^2}{1-\abar_{t-1}}\right]
=-\tfrac12\!\left[\Big(\tfrac{\alpha_t}{\beta_t}+\tfrac{1}{1-\abar_{t-1}}\Big)\|\x_{t-1}\|^2
-2\,\x_{t-1}^\top\bm b+\text{const}\right],
\]
with $\bm b=\frac{\sqrt{\alpha_t}}{\beta_t}\x_t+\frac{\sqrt{\abar_{t-1}}}{1-\abar_{t-1}}\x_0$.
Completing the square identifies a Gaussian with precision
$\tilde\beta_t^{-1}=\frac{\alpha_t}{\beta_t}+\frac{1}{1-\abar_{t-1}}=\frac{1-\abar_t}{\beta_t(1-\abar_{t-1})}$
(using $\alpha_t(1-\abar_{t-1})+\beta_t=1-\abar_t$) and mean
$\tilde{\bm\mu}_t=\tilde\beta_t\,\bm b$, which simplifies to \eqref{eq:posterior}.
\end{proof}

\subsection{Gaussian KL and the mean-matching objective}
For two Gaussians with equal covariance $\sigma_t^2\II$, the KL collapses to a
scaled squared distance of the means:
\begin{equation}
L_{t-1}=\E_q\!\left[\frac{1}{2\sigma_t^2}\big\|\tilde{\bm\mu}_t(\x_t,\x_0)-\bm\mu_\theta(\x_t,t)\big\|^2\right]+\text{const}.
\label{eq:meanmatch}
\end{equation}
So the network only needs to predict the posterior mean.

\subsection{Reparameterisation to noise prediction}
Ho et al.\ (2020) observed that predicting $\eps$ instead of $\bm\mu$ is both
simpler and empirically better. Invert \eqref{eq:kernel},
$\x_0=(\x_t-\sqrt{1-\abar_t}\,\eps)/\sqrt{\abar_t}$, and substitute into
\eqref{eq:posterior}:
\begin{equation}
\tilde{\bm\mu}_t=\frac{1}{\sqrt{\alpha_t}}\Big(\x_t-\frac{\beta_t}{\sqrt{1-\abar_t}}\,\eps\Big).
\label{eq:mu-eps}
\end{equation}
Choosing the \emph{same} functional form for the model,
$\bm\mu_\theta=\frac{1}{\sqrt{\alpha_t}}\big(\x_t-\frac{\beta_t}{\sqrt{1-\abar_t}}\,\eps_\theta(\x_t,t)\big)$,
makes the difference of means proportional to the difference of noises. Plugging
into \eqref{eq:meanmatch},
\begin{equation}
L_{t-1}=\E_{\x_0,\eps}\!\left[\frac{\beta_t^2}{2\sigma_t^2\,\alpha_t(1-\abar_t)}
\big\|\eps-\eps_\theta(\sqrt{\abar_t}\x_0+\sqrt{1-\abar_t}\eps,\,t)\big\|^2\right].
\label{eq:Lt-eps}
\end{equation}

\subsection{The simplified training objective}
Dropping the per-$t$ weight in \eqref{eq:Lt-eps} (which Ho et al.\ found improves
sample quality by up-weighting harder, higher-noise steps) yields the objective
this project optimises:
\begin{equation}
\boxed{\;
\mathcal{L}_{\text{simple}}(\theta)=
\E_{\x_0\sim p_{\text{healthy}},\;t\sim\mathcal U\{1,\dots,T\},\;\eps\sim\N(\bm0,\II)}
\Big[\big\|\eps-\eps_\theta(\sqrt{\abar_t}\,\x_0+\sqrt{1-\abar_t}\,\eps,\;t)\big\|_2^2\Big].
\;}
\label{eq:simple}
\end{equation}
The expectation over $\x_0$ is taken \emph{only over healthy patches}: the network
learns the score of the healthy data distribution. This is the exact loss in
\texttt{DiffusionProcess.training\_loss} (\texttt{src/models/diffusion.py}).

\subsection{The cosine noise schedule}
Rather than a linear $\beta_t$, we use the cosine schedule of Nichol \& Dhariwal
(2021), which avoids destroying information too quickly at the end of the chain:
\begin{equation}
\abar_t=\frac{f(t)}{f(0)},\qquad f(t)=\cos^2\!\Big(\frac{t/T+s}{1+s}\cdot\frac{\pi}{2}\Big),
\qquad \beta_t=1-\frac{\abar_t}{\abar_{t-1}},
\label{eq:cosine}
\end{equation}
with a small offset $s$ (the \texttt{squaredcos\_cap\_v2} schedule,
$\beta\in[10^{-4},2\!\times\!10^{-2}]$). It allocates more steps to the
high-SNR regime, where the fine anatomical detail that UAD depends on lives.

% ===========================================================================
\part*{Part II\quad Deterministic sampling with DDIM}
\addcontentsline{toc}{section}{Part II — DDIM}
\label{part:ddim}
% ===========================================================================

\section{Non-Markovian inference and the DDIM update}
Training fixes the marginals \eqref{eq:kernel} but \emph{not} the joint; Song et
al.\ (2021) exploit this freedom with a family of non-Markovian forward
processes sharing the same $q(\x_t\given\x_0)$, indexed by a stochasticity vector
$\sigma_t\ge0$:
\begin{equation}
q_\sigma(\x_{t-1}\given\x_t,\x_0)=\N\!\Big(\sqrt{\abar_{t-1}}\,\x_0
+\sqrt{1-\abar_{t-1}-\sigma_t^2}\cdot\frac{\x_t-\sqrt{\abar_t}\,\x_0}{\sqrt{1-\abar_t}},\;\sigma_t^2\II\Big).
\end{equation}
Because the $\eps$-objective \eqref{eq:simple} is unchanged, we reuse the trained
$\eps_\theta$. Substituting the predicted clean sample
\begin{equation}
\hat{\x}_0(\x_t,t)=\frac{\x_t-\sqrt{1-\abar_t}\,\eps_\theta(\x_t,t)}{\sqrt{\abar_t}}
\label{eq:x0hat}
\end{equation}
gives the generative update
\begin{equation}
\x_{t-1}=\sqrt{\abar_{t-1}}\,\hat{\x}_0(\x_t,t)
+\sqrt{1-\abar_{t-1}-\sigma_t^2}\;\eps_\theta(\x_t,t)+\sigma_t\,\z,\quad \z\sim\N(\bm0,\II).
\label{eq:ddim}
\end{equation}
Setting $\sigma_t=0$ for all $t$ makes the map \textbf{deterministic} (DDIM): the
only randomness in a reconstruction is the single noise vector injected at the
start. We use $\sigma_t=0$ so that pseudo-healthy reconstructions are
reproducible, and we evaluate \eqref{eq:ddim} on a short sub-sequence of
timesteps ($\le 50$ DDIM steps inside $[0,T_{\text{int}}]$), which is exact for the
marginals and $\sim\!20\times$ faster than the full chain.

% ===========================================================================
\part*{Part III\quad The PDM anomaly-detection pipeline}
\addcontentsline{toc}{section}{Part III — PDM pipeline}
\label{part:pdm}
% ===========================================================================

\section{Correlated (simplex) noise}
White Gaussian noise has a flat power spectrum, so the forward process erases
\emph{fine, pixel-scale} detail first; large coherent structures (and large
lesions) survive to high $t$. Following AnoDDPM (Wyatt et al., 2022) we instead
corrupt with \textbf{spatially correlated} noise approximating simplex noise,
built as a sum of fractal octaves
\begin{equation}
\bm n(\bm r)=\frac{1}{Z}\sum_{o=0}^{O-1} \rho^{\,o}\,g_{f_0\lambda^{o}}(\bm r),
\qquad O=6,\;\rho=0.8,\;\lambda=2,\;f_0=16,
\label{eq:simplex}
\end{equation}
where $g_{f}$ is a band of spatial frequency $f$, $\rho$ is the per-octave
\emph{persistence} (amplitude decay), $\lambda$ the \emph{lacunarity} (frequency
growth), and $Z$ normalises to unit variance. The resulting power spectrum is
concentrated at \emph{low} frequencies ($\sim 1/f$), so the corruption attacks
mid-/large-scale structure. The model is trained with this same noise
(\texttt{src/noise/simplex.py}); the objective \eqref{eq:simple} is unchanged in
form (predict the injected $\bm n$), while the schedule coefficients
$\sqrt{\abar_t},\sqrt{1-\abar_t}$ remain the mixing weights. This biases the
detector toward the blob-like spatial scale of tumours.

\section{Patch diffusion and Gaussian fusion}
\subsection{Why patches}
The model operates on overlapping $P\times P=96\times96$ patches (stride $16$)
rather than whole slices. With $\sim$196 patches per slice, a few thousand
healthy slices become $\sim$570k training samples --- a $\sim$196$\times$ data
amplification that is decisive in the low-patient pediatric regime --- and the
network is forced to learn a \emph{local} texture model of healthy tissue, which
generalises better than a global one and cannot ``memorise'' whole-slice layout.
Each patch is an independent draw of \eqref{eq:simple}.

\subsection{Fusion as confidence-weighted averaging}
A pixel $\bm u$ is covered by many patches; patch $p$ assigns it a score
$s_p(\bm u)$ that is more reliable near the patch centre (full context) than at
its edge. We therefore fuse with a centred Gaussian window
$w(\bm u)=\exp(-\|\bm u-\bm c_p\|^2/2\sigma_f^2)$, $\sigma_f=32$:
\begin{equation}
S(\bm u)=\frac{\sum_{p}w_p(\bm u)\,s_p(\bm u)}{\sum_{p}w_p(\bm u)}.
\label{eq:fusion}
\end{equation}
This is the maximum-likelihood combination of per-patch estimates under the
assumption that edge pixels carry higher variance (inverse-variance weighting),
and it removes the blocky seams a hard tiling would produce
(\texttt{GaussianPatchFuser}, \texttt{src/data/patches.py}).

\section{Partial-noise reconstruction: why anomalies ``heal''}
\label{sec:partial}
At inference we do \emph{not} start from pure noise. A test patch is encoded,
noised \emph{partially} to an intermediate level $T_{\text{int}}$ via
\eqref{eq:kernel}, then denoised deterministically with \eqref{eq:ddim} back to
$t=0$:
\begin{equation}
\x_{T_{\text{int}}}=\sqrt{\abar_{T_{\text{int}}}}\,\x_0+\sqrt{1-\abar_{T_{\text{int}}}}\,\bm n
\;\xrightarrow[\text{DDIM},\,\sigma=0]{\;\eps_\theta\;}\;\hat{\x}_0^{\,\text{healthy}}.
\end{equation}
\textbf{Heuristic projection argument.} Let $\mathcal M_h$ be the manifold of
healthy patches (the support of $p_{\text{healthy}}$ the model learned). Decompose
an anomalous patch as $\x_0=\x_h+\bm a$, $\x_h\in\mathcal M_h$, $\bm a$ the lesion
perturbation. Forward noising adds isotropic-in-law corruption of magnitude
$\sqrt{1-\abar_{T_{\text{int}}}}$; if the lesion's energy lies within the band the
corruption destroys, then $\x_{T_{\text{int}}}$ is, with high probability, closer to
the noised \emph{healthy} sample than to any noised anomalous one. The reverse
chain --- trained only on $\mathcal M_h$ --- then denoises toward $\mathcal M_h$,
returning $\hat{\x}_0^{\text{healthy}}\approx\x_h$ and \emph{erasing} $\bm a$.
Healthy tissue ($\bm a=\bm 0$) is reconstructed faithfully. The residual
\begin{equation}
\bm r=\big|\x_0-\hat{\x}_0^{\text{healthy}}\big|\approx|\bm a|
\end{equation}
therefore localises the anomaly. The choice of $T_{\text{int}}$ is the central
trade-off: too small fails to erase large lesions (the band is too narrow), too
large discards patient-specific healthy anatomy (degrading the residual
baseline). This motivates multi-scale scoring.

\section{Multi-scale, modality-weighted anomaly score}
\subsection{Per-modality residual and the CE channel}
Modalities differ in clinical salience, so the scalar residual is a weighted sum
over channels, augmented with an explicit contrast-enhancement channel
$\mathrm{CE}=t1c-t1n$ (a marker of active tumour invisible to pipelines that drop
native T1):
\begin{equation}
r(\bm u)=\sum_{k\in\mathcal K}\bar w_k\,\big|\x_0^{(k)}(\bm u)-\hat{\x}_0^{(k)}(\bm u)\big|,
\quad
r_{\mathrm{CE}}=\big|(\x_0^{t1c}-\x_0^{t1n})-(\hat{\x}_0^{t1c}-\hat{\x}_0^{t1n})\big|,
\label{eq:residual}
\end{equation}
with weights $\bm w\propto(0.5,1.0,0.75,1.5)$ for $(t1n,t1c,t2w,t2f)$ and
$w_{\mathrm{CE}}=1.0$, normalised so $\sum_k\bar w_k=1$
(\texttt{weighted\_residual}, \texttt{src/scoring/residual.py}). FLAIR ($t2f$) and
contrast dominate, mirroring radiological reasoning.

\subsection{Scale fusion}
We reconstruct at several noise levels $T\in\mathcal T=\{50,150,250\}$ ---
low $T$ for fine/texture lesions (high SNR retains detail), high $T$ for large
structural lesions --- producing a fused per-scale map $S_T$ via
\eqref{eq:fusion}, then combine them with fixed weights:
\begin{equation}
\boxed{\;S(\bm u)=\sum_{T\in\mathcal T}\omega_T\,S_T(\bm u)\cdot m_{\text{brain}}(\bm u),
\qquad \bm\omega=(0.25,0.45,0.30).\;}
\label{eq:multiscale}
\end{equation}
$m_{\text{brain}}$ is the brain mask (a 2-px erosion removes the skull rim while
sparing cortical lesions). Equation~\eqref{eq:multiscale} is computed by
\texttt{MultiScaleScorer.score\_slice}.

\section{Healthy-set calibration and thresholding}
The operating threshold is set \emph{without ever seeing a lesion}. On $N$
held-out healthy slices we pool all brain-voxel scores and take a high
percentile,
\begin{equation}
\tau=\mathrm{Percentile}_{95}\big(\{S(\bm u):\bm u\in m_{\text{brain}},\;\x\text{ healthy}\}\big),
\qquad \hat P(\bm u)=\mathbf 1\!\left[S(\bm u)>\tau\right].
\label{eq:tau}
\end{equation}
By construction $\tau$ passes (in expectation) only the top $5\%$ of healthy
voxels, i.e.\ a controlled, label-free false-positive rate $\alpha=0.05$. A
percentile is a Wilks (1941) order statistic: $\sim20/\alpha=400$ pooled samples
suffice for a stable estimate, comfortably met by $N=150$--$200$ slices
($\sim10^6$ voxels). This is \texttt{src/calibration/calibrator.py}.

\section{Evaluation metrics}
Given continuous scores $S$ and binary ground truth $G$ over brain voxels:
\begin{align}
\mathrm{AUROC}&=\Pr\!\big(S(\bm u_+)>S(\bm u_-)\big)
=\int_0^1 \mathrm{TPR}\,d\,\mathrm{FPR}, \\[2pt]
\mathrm{AUPRC}&=\int_0^1 \mathrm{Precision}\,d\,\mathrm{Recall}, \\[2pt]
\mathrm{Dice}(\hat P,G)&=\frac{2\,|\hat P\cap G|}{|\hat P|+|G|}
=\frac{2\,\mathrm{TP}}{2\,\mathrm{TP}+\mathrm{FP}+\mathrm{FN}}\;(=F_1).
\label{eq:metrics}
\end{align}
AUROC/AUPRC are \emph{threshold-free} (pure ranking quality); Dice depends on the
threshold and is reported at three operating points --- the calibrated $\tau$ of
\eqref{eq:tau} (deployable), the best single global threshold (realistic
ceiling), and the per-slice oracle (upper bound). The spread between them
diagnoses whether error is in the \emph{score} or only in the \emph{threshold}.

% ===========================================================================
\part*{Part IV\quad Explainability (XAI) mathematics}
\addcontentsline{toc}{section}{Part IV — XAI}
\label{part:xai}
% ===========================================================================

\section{Counterfactual explanation (primary)}
A diffusion UAD model answers an inherently \emph{counterfactual} question:
``what would this brain look like if it were healthy?'' The reconstruction
$\hat{\x}_0^{\text{healthy}}$ \emph{is} the counterfactual $\x'$, and the anomaly
map $|\x-\x'|$ is the counterfactual difference --- the minimal change that moves
the input onto the healthy manifold (Wachter et al., 2017; Sanchez et al., 2022).
This explanation is \textbf{faithful by construction}: it is exactly the quantity
the model thresholds, not a post-hoc surrogate. The DDIM map of
\S\ref{sec:partial} approximates the projection onto $\mathcal M_h$ in the
metric induced by the learned score, so $\x'$ is, approximately, the nearest
healthy explanation of $\x$. The ``healing trajectory'' visualises the iterates
of \eqref{eq:ddim} in-painting the lesion.

\section{Grad-CAM for a denoiser (model-focus)}
Counterfactuals say \emph{what} changed; Grad-CAM (Selvaraju et al., 2017) says
\emph{where the network looked}. A denoiser has no class logit, so we explain its
own training target evaluated at the centre scoring scale,
\begin{equation}
y(\x_0,t)=\big\|\eps_\theta(\x_t,t)-\bm n\big\|_2^2,
\qquad \x_t=\sqrt{\abar_t}\,\x_0+\sqrt{1-\abar_t}\,\bm n,
\label{eq:gradcam-target}
\end{equation}
which is large precisely where the network \emph{fails} to denoise as if healthy
--- the anomaly. Let $\bm A^{k}$ be the $k$-th channel of a high-resolution
decoder feature map (the last up-block). The Grad-CAM weight and map are
\begin{equation}
\alpha_k=\frac{1}{HW}\sum_{i,j}\frac{\partial y}{\partial A^{k}_{ij}},
\qquad
L_{\text{Grad-CAM}}=\ReLU\!\Big(\sum_k \alpha_k\,\bm A^{k}\Big),
\label{eq:gradcam}
\end{equation}
upsampled to $P\times P$, min--max normalised, and fused across patches with
\eqref{eq:fusion} into a slice-level model-focus map directly comparable to $S$
and $G$. The gradients in \eqref{eq:gradcam} are obtained by one backward pass of
$y$ onto the cached activation (\texttt{src/xai/gradcam.py}); this is the only
part of the pipeline that runs \emph{with} gradients.

\section{Additive attribution: modality and scale}
Because both the residual \eqref{eq:residual} and the multi-scale score
\eqref{eq:multiscale} are \emph{linear} combinations, attribution is exact (no
approximation): each term is a literal contribution.
\begin{equation}
S=\sum_{k\in\mathcal K}\underbrace{\bar w_k\, r_k}_{\text{modality } k}
=\sum_{T\in\mathcal T}\underbrace{\omega_T\,S_T}_{\text{scale } T}.
\end{equation}
The \emph{dominant modality} $\arg\max_k \sum_{\bm u}\bar w_k r_k(\bm u)$ indicates
whether a lesion is driven by FLAIR (edema) or contrast (active tumour); the
\emph{dominant scale} indicates lesion size (low-$T$ = fine, high-$T$ =
structural). These are \texttt{modality\_attribution} / \texttt{scale\_attribution}
(\texttt{src/xai/attribution.py}). An attention-rollout view (Abnar \& Zuidema,
2020), aggregating the UNet self-attention magnitudes, is available as a
complementary model-internal saliency.

% ===========================================================================
\part*{Part V\quad Reading the numbers}
\addcontentsline{toc}{section}{Part V — Reading the numbers}
\label{part:numbers}
% ===========================================================================

\section{Training numbers}
\paragraph{The loss scale.} The injected noise $\bm n$ is normalised to unit
variance, so a trivial predictor $\eps_\theta\!\equiv\!\bm0$ attains
$\mathcal L_{\text{simple}}=\E\|\bm n\|^2/(CP^2)\approx 1.0$. Hence
$\mathcal L_{\text{simple}}\in[0,1]$ in practice and
$1-\mathcal L_{\text{simple}}$ is, approximately, the \emph{fraction of noise
variance the model explains}. A training loss of $\sim$0.09 corresponds to
explaining $\sim$90\% of the injected-noise variance --- the network has learned
a strong healthy-denoising score function.
\paragraph{Train vs.\ validation.} Validation runs the EMA weights on held-out
healthy patches. A validation loss \emph{higher} than training (e.g.\ $0.26$ vs.\
$0.09$ early on) is expected --- different patches, EMA lag --- and is not
over-fitting; what matters is that \emph{validation} trends \emph{down} across
epochs and then plateaus, which is the early-stopping signal (patience $6$). A
validation loss that rises while training falls would indicate over-fitting; the
patch amplification of \S\ref{part:pdm} makes this unlikely.
\paragraph{Steps, not epochs.} With the 250k-patch-per-epoch cap, one epoch is
$\sim$977 optimiser steps; $\sim$4--5 capped epochs already exceed the gradient
budget of a far longer un-capped run on a smaller dataset, which is why a
few-epoch pilot is legitimate. Convergence should be read off the
\emph{validation curve}, not the epoch count.

\section{Inference numbers}
\paragraph{AUROC vs.\ Dice.} A high AUROC with a low calibrated Dice means the
score \emph{ranks} lesion voxels correctly but the single global threshold
\eqref{eq:tau} cannot cleanly separate them --- typically because a second,
non-lesion structure (e.g.\ a reconstruction edge rim) also exceeds $\tau$. The
\emph{best-global} and \emph{oracle} Dice quantify how much of the gap is
threshold-induced versus intrinsic to the score: a large oracle$-$calibrated gap
means ``fix the threshold/artefact,'' a small one means ``improve the score.''
\paragraph{Per-scale / per-modality readouts.} The scale weights $\bm\omega$ and
the dominant-modality statistic turn each detection into a clinically legible
statement (lesion size; FLAIR- vs.\ contrast-driven), and the per-scale Dice
isolates which $T$ contributes most --- guidance for tuning $\mathcal T$.
\paragraph{Grad-CAM sanity.} A trustworthy model-focus map \eqref{eq:gradcam}
concentrates on the lesion and co-localises with both the anomaly score $S$ and
the ground truth $G$; diffuse or edge-locked Grad-CAM indicates the network is
keying on artefacts rather than pathology, an independent check on the score.

\bigskip
\hrule
\bigskip
\noindent\textit{Code cross-reference.} Eq.~\eqref{eq:simple}: \texttt{models/diffusion.py};
Eqs.~\eqref{eq:x0hat}--\eqref{eq:ddim}: \texttt{models/diffusion.py} (DDIM);
Eqs.~\eqref{eq:simplex}: \texttt{noise/simplex.py};
Eqs.~\eqref{eq:fusion},\eqref{eq:multiscale},\eqref{eq:residual}: \texttt{scoring/};
Eq.~\eqref{eq:tau}: \texttt{calibration/calibrator.py};
Eqs.~\eqref{eq:gradcam-target}--\eqref{eq:gradcam}: \texttt{xai/gradcam.py}.

\section*{Key references}
\addcontentsline{toc}{section}{Key references}
\begin{enumerate}[label={[\arabic*]},leftmargin=2.4em]
\item Sohl-Dickstein, Weiss, Maheswaranathan, Ganguli (2015). \textit{Deep Unsupervised Learning using Nonequilibrium Thermodynamics}. ICML.
\item Ho, Jain, Abbeel (2020). \textit{Denoising Diffusion Probabilistic Models}. NeurIPS.
\item Song, Meng, Ermon (2021). \textit{Denoising Diffusion Implicit Models}. ICLR.
\item Nichol, Dhariwal (2021). \textit{Improved Denoising Diffusion Probabilistic Models}. ICML.
\item Wyatt, Leach, Schmon, Willcocks (2022). \textit{AnoDDPM: Anomaly Detection with DDPMs using Simplex Noise}. CVPR-W.
\item Behrendt, Bhattacharya, Krüger, Opfer, Schlaefer (2023). \textit{Patched Diffusion Models for UAD in Brain MRI}. MIDL.
\item Bercea et al.\ (2023). \textit{Generalizing Unsupervised Anomaly Detection}. MICCAI.
\item Baur, Denner, Wiestler, Navab, Albarqouni (2021). \textit{Autoencoders for Unsupervised Anomaly Segmentation in Brain MR}. Med.\ Image Anal.
\item Selvaraju et al.\ (2017). \textit{Grad-CAM}. ICCV.
\item Wachter, Mittelstadt, Russell (2017). \textit{Counterfactual Explanations without Opening the Black Box}. Harvard JOLT.
\item Sanchez, Kascenas, Liu, O'Neil, Tsaftaris (2022). \textit{Healthy/Pathological Brain Counterfactuals with Diffusion Models}. arXiv:2203.08089.
\item Abnar, Zuidema (2020). \textit{Quantifying Attention Flow in Transformers}. ACL.
\item Wilks (1941). \textit{Determination of Sample Sizes for Setting Tolerance Limits}. Ann.\ Math.\ Stat.
\end{enumerate}

\end{document}
